\section{Einleitung}
\label{sec:einleitung}

Dieses Kapitel führt in die Ausgangssituation der Arbeit ein und leitet aus den
Herausforderungen der Brownfield-Transformation die zugrunde liegende Problemstellung
und Zielsetzung ab. Anschließend werden der Untersuchungsgegenstand abgegrenzt und
der Aufbau der Arbeit erläutert.

\subsection{Ausgangssituation und Motivation}
\label{sec:motivation}

Model-Based Systems Engineering (MBSE) verfolgt das Ziel, systembezogene
Engineering-Information in formalisierten und miteinander verknüpften Modellen
abzubilden und dadurch deren Nutzung über unterschiedliche
Entwicklungsaktivitäten hinweg zu unterstützen. Insbesondere bei komplexen
technischen Systemen können dadurch Zusammenhänge nachvollziehbarer
dargestellt, Engineering-Artefakte konsistenter miteinander verbunden und
enthaltene Informationen stärker maschinell verarbeitet werden
\cite{Hendriks2022,Campo2023}.

Die Einführung von MBSE ist mit erheblichem organisatorischem und technischem
Aufwand verbunden. Neben der Auswahl und Einführung geeigneter Werkzeuge
müssen Modellierungsprozesse etabliert, Kompetenzen aufgebaut und bestehende
Engineering-Information in modellbasierte Strukturen überführt werden
\cite{Call2022,Henderson2024}. Ein wesentlicher Teil dieses Aufwands fällt
bereits in frühen Phasen einer MBSE-Einführung an. Bevor die Vorteile einer
modellbasierten Arbeitsweise genutzt werden können, müssen
Modellierungsumgebungen, Arbeitsrichtlinien und initiale Modellinhalte
aufgebaut werden. Dies kann die Einführung insbesondere dann erschweren,
wenn dafür nur begrenzte personelle und finanzielle Ressourcen zur Verfügung
stehen \cite{Call2022}. Für kleine und mittlere Unternehmen stellen neben
fehlender Erfahrung insbesondere die wirtschaftliche Rechtfertigung,
Schulungsbedarf und verfügbare Unterstützung relevante Herausforderungen dar
\cite{Nanfuka2023}.

Diese Ausgangslage wird dadurch geprägt, dass Unternehmen eine Transformation
zu MBSE typischerweise mit bereits gewachsenen Informationsbeständen
beginnen. Technische Dokumente, Tabellen, Spezifikationen und weitere
Systembeschreibungen enthalten Engineering-Information, die über längere
Zeiträume und in unterschiedlichen Strukturen entstanden ist. Eine
Brownfield-Transformation umfasst deshalb sowohl den Aufbau modellbasierter
Strukturen als auch den Umgang mit diesen bestehenden Informationsbeständen.
Erfahrungen aus der industriellen MBSE-Einführung zeigen, dass
Legacy-Artefakte und modellbasierte Arbeitsweisen während einer solchen
Transformation über längere Zeit parallel bestehen können
\cite{Henderson2024}.

Auch die in der Literatur beschriebenen wirtschaftlichen und
organisatorischen Effekte von MBSE ergeben hinsichtlich Aufwand und Nutzen
kein einheitliches Bild. Vorteilen hinsichtlich Entwicklungszeit, Qualität
und Zusammenarbeit stehen Zeit-, Kosten- und Arbeitsaufwände für Einführung
und Anwendung gegenüber; zugleich ist die empirische Absicherung vieler
dieser Aussagen begrenzt \cite{Campo2023}. Damit gewinnt insbesondere die
Frage an Bedeutung, wie wiederkehrende Aufwände bei der Überführung
bestehender Engineering-Information unterstützt werden können.

Der praktische Ausgangspunkt dieser Arbeit liegt in der schrittweisen
Modellierung weiterer Systembereiche in einem industriellen Kontext. Dabei
entsteht wiederholt Aufwand für die Erfassung zusätzlicher
Engineering-Information, ihre fachliche Einordnung und ihre Abstimmung mit
bereits bestehenden Modellinhalten. Diese Beobachtung motiviert die Suche
nach Verfahren, mit denen bestehende Engineering-Information wiederholbar und
nachvollziehbar in formale Modelle überführt werden kann.

Generative Künstliche Intelligenz (KI) eröffnet hierfür die Möglichkeit,
auch schwach strukturierte und natürlichsprachliche Informationen
automatisiert zu verarbeiten. Für den Einsatz im Engineering muss diese
Fähigkeit durch Mechanismen ergänzt werden, die Herkunft, Interpretation und
fachliche Entscheidung nachvollziehbar erhalten. Daraus ergibt sich die
Frage, wie eine KI-Unterstützung gestaltet werden kann, mit der bestehende
Engineering-Information unter Einbindung menschlicher
Entscheidungsbefugnis kontrolliert in MBSE-Modelle überführt und dort
weiterverwendet werden kann \cite{SheardPickard2025}. Diese Fragestellung
bildet den Ausgangspunkt der vorliegenden Arbeit.


\subsection{Problemstellung}
\label{sec:problemstellung}

Die Überführung bestehender Engineering-Information in eine modellbasierte
Repräsentation umfasst mehr als eine Formatkonvertierung. Brownfield-Projekte
verfügen über gewachsene Informationsbestände in Form von Spezifikationen,
Tabellen, technischen Dokumenten und weiteren Engineering-Artefakten. Diese
unterscheiden sich hinsichtlich Struktur, Granularität, Terminologie und
Explizitheit und müssen daher fachlich interpretiert und in einen geeigneten
Modellkontext eingeordnet werden \cite{Bressan.2000}. Erfahrungen aus der
industriellen MBSE-Einführung zeigen zugleich, dass bestehende
Engineering-Artefakte während einer Transformation parallel zu den
entstehenden modellbasierten Strukturen weiterverwendet werden können
\cite{Henderson2024}.

Bleiben solche Engineering-Artefakte Bestandteil regulärer Entwicklungs- und
Änderungsprozesse, können spätere Änderungen eine erneute Betrachtung und
gegebenenfalls Anpassung bereits erstellter Modellinformationen erfordern.
Die Brownfield-Transformation kann sich damit über mehrere sukzessive
Transformationsschritte erstrecken. Mit der Einbeziehung weiterer
Systembereiche entsteht erneut Aufwand für Erfassung, fachliche Einordnung,
Modellbildung, Integration und die Abstimmung mit bereits bestehenden
Teilmodellen.

Dieser wiederkehrende Modellierungs- und Integrationsaufwand wurde auch im
Rahmen der industriellen Themenfindung dieser Arbeit wahrgenommen.
Abbildung~\ref{fig:problemwahrnehmung_transformation} stellt diesen
Ausgangspunkt einem angestrebten KI-unterstützten Transformationspfad
schematisch gegenüber.

\begin{figure}[htbp]
    \centering
    \includegraphics[width=0.9\textwidth]
      {Graphiken/1.2_Problemdarstellung.png}
    \caption[Brownfield-Transformation und angestrebter Transformationspfad]{
        Schematische Darstellung der im Rahmen der industriellen
        Themenfindung wahrgenommenen wiederkehrenden Modellierungs- und
        Integrationsaufwände bei sukzessiver Brownfield-Transformation sowie
        eines angestrebten KI-unterstützten Transformationspfads
        (eigene Darstellung)}
    \label{fig:problemwahrnehmung_transformation}
\end{figure}

Die lila Kurve beschreibt schematisch den wiederkehrenden manuellen Aufwand
einer sukzessiven Brownfield-Transformation. Mit der Erschließung weiterer
Systembereiche entstehen erneut Aufwände für die Erfassung und fachliche
Einordnung bestehender Engineering-Information sowie für Modellbildung,
Integration und die Abstimmung mit bereits vorhandenen Modellinhalten.

Die grüne Kurve visualisiert die motivierende Zielvorstellung eines
KI-unterstützten Transformationspfads. Der Aufbau und die Konfiguration eines
unterstützenden Transformationssystems verursachen zunächst zusätzlichen
Aufwand. Bei der wiederholten Einbindung weiterer Engineering-Information
sollen wiederverwendbare und soweit geeignet automatisierte
Verarbeitungsschritte dazu beitragen, den erneut entstehenden manuellen
Aufwand zu reduzieren. Die vorliegende Arbeit greift dieses Zielbild auf und
untersucht die architektonischen Voraussetzungen für ein entsprechendes
Transformationssystem.

KI-generierte Inhalte können syntaktisch und semantisch plausibel erscheinen
und zugleich unvollständige oder fachlich fehlerhafte Aussagen enthalten
\cite{Topcu2025,Barke2023,Sergeyuk2025}. Die Verlässlichkeit eines
KI-unterstützten Transformationsprozesses erfordert daher zusätzlich zur
Erzeugung plausibler Modellinhalte eine kontrollierte Herleitung der darin
repräsentierten Engineering-Information.

Hierzu muss nachvollziehbar bleiben, aus welchen Quellen eine Information
stammt, wie sie interpretiert wurde, an welchen Stellen Unsicherheit oder
Mehrdeutigkeit besteht und durch welche fachlichen Entscheidungen sie für
die weitere Verarbeitung autorisiert wurde. Die nachvollziehbare Herkunft
und Entstehung einer Information wird im weiteren Verlauf dieser Arbeit als
\emph{Provenienz} bezeichnet. Ergänzend muss festgelegt werden, wie fachlich
autorisierte Information in eine geeignete Modellstruktur eingeordnet und in
der vorgesehenen Zielnotation repräsentiert wird.

Die Problemstellung dieser Arbeit besteht damit in der Gestaltung eines
Transformationsprozesses, der heterogene Engineering-Bestandsinformation,
probabilistische KI-Unterstützung, Provenienz, menschliche
Entscheidungsbefugnis und die Überführung in SysML~v2 in einer
zusammenhängenden und nachvollziehbaren Informationsverarbeitung verbindet.


\subsection{Zielsetzung der Arbeit}
\label{sec:zielsetzung}

Ziel dieser Arbeit ist die Entwicklung und Untersuchung eines
Architekturkonzepts für die KI-gestützte Transformation bestehender
heterogener Engineering-Information in SysML~v2. Im Mittelpunkt steht die
Frage, wie ein kontrollierter und nachvollziehbarer Transformationsprozess
zwischen bestehenden Informationsbeständen und einer formalen
Zielrepräsentation gestaltet werden kann.

Das Architekturkonzept berücksichtigt hierzu die Provenienz verarbeiteter
Informationen, die kontrollierte Nutzung probabilistischer
KI-Interpretationen, die gezielte Einbindung menschlicher Entscheidungen
sowie technisch überprüfbare nachgelagerte Verarbeitungsschritte. Dabei
werden probabilistische Interpretation, fachliche Entscheidungsbefugnis und
Modellbildung als unterschiedliche Aufgaben innerhalb des
Transformationsprozesses behandelt. Die konkrete Ausgestaltung dieser
Aufgaben und ihrer Übergänge wird in Kapitel~\ref{sec:architecture}
systematisch hergeleitet.

Die Architektur wird unabhängig von einer konkreten
Softwareimplementierung beschrieben. Zur Untersuchung ihrer zentralen
Annahmen wurde im Rahmen dieser Arbeit ein Prototyp des
Transformationssystems entwickelt. Dieser realisiert ausgewählte
Informationsverarbeitungs-, Kontroll- und Übergangsmechanismen in einem
durchgängigen Verarbeitungspfad von registrierten Engineering-Quellen bis zu
einer SysML-v2-Repräsentation. Die prototypische Realisierung dient damit als
Untersuchungsinstrument für das entwickelte Architekturkonzept.

Die Untersuchung erfolgt exemplarisch anhand kontrolliert aufbereiteter
Engineering-Information aus einem industriell geprägten Anwendungskontext.
Dabei wird insbesondere betrachtet, an welchen Stellen probabilistische
Verarbeitung eingesetzt werden kann, welche Informationen für
nachgelagerte Schritte autorisiert werden und an welchen Übergängen
menschliche Entscheidungen erforderlich sind. Die erzeugte
SysML-v2-Repräsentation wird technisch validiert und der resultierende
Modellzustand anschließend menschlich geprüft und freigegeben.

Die schrittweise Einbeziehung weiterer Engineering-Artefakte und die
fortlaufende Veränderung bestehender Informationsbestände werden als
relevante Eigenschaften einer Brownfield-Transformation berücksichtigt.
Stabile Identitäten, Provenienz sowie Beziehungen zwischen aufeinander
folgenden Informations- und Modellzuständen schaffen dabei eine Grundlage,
auf der weiterführende Mechanismen zur inkrementellen Aktualisierung
aufbauen können.

Die Arbeit untersucht damit die architektonischen Bedingungen, unter denen
KI zur kontrollierten Überführung bestehender Engineering-Information in
SysML~v2 eingesetzt werden kann. Die mögliche Verringerung wiederkehrender
Modellierungs- und Integrationsaufwände bildet die praktische Motivation
dieses Ansatzes; die wissenschaftliche Untersuchung konzentriert sich auf
die hierfür erforderliche Informationsverarbeitungsarchitektur und ihre
prototypische Umsetzung.

\subsection{Abgrenzung}
\label{sec:abgrenzung}

Die prototypische Realisierung dient als Untersuchungsinstrument für die in
dieser Arbeit entwickelten Architekturmechanismen. Der Untersuchungsumfang
konzentriert sich auf die kontrollierte Verarbeitung bestehender
Engineering-Information, ihre fachliche und modellbezogene Autorisierung,
die Überführung in eine SysML-v2-Repräsentation sowie die dabei erforderliche
Provenienz und Nachverfolgbarkeit.

Menschliche Entscheidungen sind dabei ein vorgesehener Bestandteil des
Transformationsprozesses. Sie werden insbesondere dort eingesetzt, wo
Engineering-Information fachlich bewertet und autorisiert, im Zielmodell
eingeordnet oder der erzeugte Modellzustand geprüft und freigegeben wird.
Die Arbeit untersucht damit eine kontrollierte KI-Unterstützung mit
expliziter menschlicher Entscheidungsbefugnis.

Die fortlaufende Veränderung bestehender Engineering-Informationsbestände
wird als relevante Eigenschaft einer Brownfield-Transformation
berücksichtigt. Die entwickelte Architektur erhält hierzu Identitäten,
Provenienz und Beziehungen zwischen aufeinanderfolgenden Informations- und
Modellzuständen. Eine kontinuierliche Synchronisation bestehender
Engineering-Artefakte mit dem erzeugten Systemmodell sowie eine vollständige
Analyse der Auswirkungen von Änderungen auf das Modell sind nicht Bestandteil
der prototypischen Untersuchung. Die hierfür erforderlichen weiterführenden
Mechanismen werden im Ausblick aufgegriffen.

Die prototypische Untersuchung verwendet exemplarisch kontrolliert
aufbereitete Engineering-Information aus einem bestehenden industriellen
Systemkontext der PreciPoint GmbH. Gegenstand der Evaluation sind die
entwickelten Architekturmechanismen und ihr Zusammenspiel im
Transformationsprozess. Die Untersuchung liefert damit Aussagen über ihre
prototypische Umsetzbarkeit und die dabei beobachteten
Informations-, Entscheidungs- und Modellübergänge.

Wirtschaftlichkeit, Produktivitätssteigerung und Zeitersparnis gegenüber
einer manuellen Modellierung werden als mögliche Motivation einer
KI-gestützten Brownfield-Transformation betrachtet, jedoch nicht quantitativ
untersucht. Die technische Validierung der erzeugten
SysML-v2-Repräsentation adressiert deren Verarbeitbarkeit in der verwendeten
SysML-v2-Toolumgebung. Die fachliche Bewertung erfolgt davon getrennt durch
die im Transformationsprozess vorgesehenen menschlichen Entscheidungen.

\subsection{Aufbau der Arbeit}
\label{sec:aufbau}

Kapitel~2 ordnet die Problemstellung in den wissenschaftlichen Stand ein.
Dazu werden die für diese Arbeit relevanten Grundlagen zu MBSE und
SysML~v2, zur Transformation bestehender Engineering-Information sowie zu
KI-gestützter Informationsverarbeitung und den hierfür relevanten
Kontrollmechanismen betrachtet. Aus den dabei identifizierten offenen
Fragestellungen werden die Forschungsfragen der Arbeit abgeleitet.

Kapitel~3 beschreibt das methodische Vorgehen zur Bearbeitung dieser
Forschungsfragen. Dazu gehören die Literaturrecherche, die modellbasierte
Architekturentwicklung, die methodischen Gestaltungsprinzipien der
KI-gestützten Informationsverarbeitung sowie die prototypische Untersuchung
und Evaluation des entwickelten Architekturkonzepts.

Darauf aufbauend entwickelt Kapitel~4 die konzeptionelle
Informationsverarbeitungsarchitektur für die kontrollierte Transformation
bestehender Engineering-Information in SysML~v2. Die dort eingeführten
Begriffe beschreiben die fachlichen Informations-, Entscheidungs- und
Modellbildungsmechanismen der Architektur.

Kapitel~5 beschreibt deren prototypische Realisierung. Die konzeptionellen
Architekturmechanismen aus Kapitel~4 werden dabei auf das entwickelte
SysML-v2-Architekturmodell und die ausführbare Implementierung bezogen.
Soweit konkrete Elemente des Architekturmodells beschrieben werden,
orientieren sich die verwendeten Bezeichnungen an den Bezeichnungen des
Modells. Dadurch können die textuelle Architekturbeschreibung und das
zugehörige SysML-v2-Modell unmittelbar miteinander abgeglichen werden.

Kapitel~6 beschreibt das Vorgehen zur Verifikation und Validierung der
prototypischen Realisierung. Die dabei beobachteten Ergebnisse werden davon
getrennt in Kapitel~7 dargestellt. Kapitel~8 ordnet diese Ergebnisse in das
entwickelte Architekturkonzept und den wissenschaftlichen Kontext ein und
führt sie zur Beantwortung der Forschungsfragen zusammen.

Kapitel~9 fasst die wesentlichen Erkenntnisse der Arbeit zusammen und zeigt
darauf aufbauend weiterführende Entwicklungsmöglichkeiten und
Forschungsrichtungen auf.